% `tex.web` §216 and §663: which line a paragraph *began* on.
%
% §1091's `new_graf` calls `push_nest`, and §216 stamps the new level with
% `mode_line := line`. §663's `common_ending` reports a bad line against it:
%
%     print(") in paragraph at lines "); print_int(mode_line); …
%
% Every nest level has its own `mode_line`, so a box opened in the middle of a
% paragraph gets one and the paragraph keeps its own. We had a single field,
% used both as the paragraph's line and as the slot the *next* box is stamped
% from — so any `\hbox` between the start of a paragraph and its end
% overwrote it, and the warning named the box's line as the paragraph's start.
%
% That is the one number a reader uses to find the offending paragraph, and a
% wrong one sends them to a line that is not even part of it.
\catcode`\{=1 \catcode`\}=2
\font\r=cmr10 \r \hsize=100pt \parindent=0pt \hbadness=100 \vbadness=100
\tolerance=10000

% A paragraph that opens a box in the middle of itself. Started on the first
% line below, ended on the last.
A\hskip 0pt plus 60pt B
\setbox0\hbox{Q}
\setbox1\vbox{\hbox{Q}}
C\par

% ...and one that opens none, so the two answers coincide and the case still
% says which is which.
A\hskip 0pt plus 60pt B\par

% A box *before* a paragraph must not be its line either.
\setbox2\hbox{Q}
A\hskip 0pt plus 60pt B\par

% A paragraph inside a \vbox, with a box inside it.
\setbox3\vbox{A\hskip 0pt plus 60pt B
  \setbox4\hbox{Q}
  C\par}

% \noindent starts one the same way, and so does \indent.
\noindent A\hskip 0pt plus 60pt B
\setbox5\hbox{Q}
C\par
\indent A\hskip 0pt plus 60pt B
\setbox6\hbox{Q}
C\par
\end
